\documentclass[../../main.tex]{subfiles}% 注意这里的文件路径不能用 ./main.tex ,否则用latexmk编译子文件会报错
\graphicspath{{\subfix{./image/}}{\subfix{../../image/}}} % 指定图片引用路径,后续可以直接使用图片文件名
% 如果图片存放在上级目录,则使用 ../../image/ 作为备用路径

% 例如:
% \begin{figure}[H]
% \centering
% \includegraphics[width=0.3\textwidth]{图.png}
% \caption{}
% \label{figure:图}
% \end{figure}
% 注意:上述\label{}一定要放在\caption{}之后,否则引用图片序号会只会显示??.

\begin{document}

% SOURCE p=55; printed=40; object=ch01-sec06-title; region=middle
\section{注记与参考文献}\label{section:ch01-notes-references}

% SOURCE p=55; printed=40; object=ch01-history-intro; region=middle
测度论的历史与积分理论的历史密切相关，关于后者的评论将在\cref{section:ch02-polar-coordinates}中给出。

% SOURCE p=55; printed=40; object=ch01-notes-banach-tarski; region=middle-lower
\cref{section:ch01-introduction}：Banach–Tarski 悖论最早出现于文献\cite{folland-bib-011}，但下面这个变体可追溯到 Hausdorff\cite{folland-bib-068}：
\begin{proposition}[Hausdorff 的球面分解]\label{proposition:ch01-hausdorff-sphere-decomposition}
$\mathbb{R}^3$ 中的单位球面 $\{x\in\mathbb{R}^3:|x|=1\}$ 是四个集合 $E_1,\ldots,E_4$ 的不交并，满足：
\begin{enumerate}
\item\label{item:ch01-hausdorff-paradox-countable} $E_1$ 可数；
\item\label{item:ch01-hausdorff-paradox-rotations} 集合 $E_2$、$E_3$、$E_4$ 及 $E_3\cup E_4$ 都可以通过旋转互相得到。
\end{enumerate}
\end{proposition}
Banach–Tarski 悖论及 Hausdorff 的结果的初等论述可见 Stromberg\cite{folland-bib-146}。

% SOURCE p=55; printed=40; object=ch01-notes-transfinite-sigma-generation; region=lower
\cref{section:ch01-sigma-algebras}：我们对集合族 $\mathscr{E}\subseteq\mathscr{P}(X)$ 所生成的 $\sigma$-代数 $\mathscr{M}(\mathscr{E})$ 的刻画是非构造性的，因此可能会问如何从 $\mathscr{E}$ 显式地得到 $\mathscr{M}(\mathscr{E})$。答案相当复杂。可以从如下过程开始：令 $\mathscr{E}_1=\mathscr{E}\cup\{E^{\mathrm{c}}:E\in\mathscr{E}\}$；当 $j>1$ 时，定义 $\mathscr{E}_j$ 为所有 $\mathscr{E}_{j-1}$ 中集合的可数并及这类并集的补集组成的族。令 $\mathscr{E}_\omega=\bigcup_1^\infty\mathscr{E}_j$：是否有 $\mathscr{E}_\omega=\mathscr{M}(\mathscr{E})$？一般来说，没有。$\mathscr{E}_\omega$ 对取补集封闭，但若对每个 $j$ 有 $E_j\in\mathscr{E}_j\setminus\mathscr{E}_{j-1}$，则没有理由认为 $\bigcup_1^\infty E_j$ 属于 $\mathscr{E}_\omega$。必须重新开始。更确切地说，必须通过超限归纳法对每个可数序数 $\alpha$ 定义 $\mathscr{E}_\alpha$：若 $\alpha$ 有直接前驱 $\beta$，则 $\mathscr{E}_\alpha$ 为所有 $\mathscr{E}_\beta$ 中集合的可数并及这类并集的补集组成的族；否则，$\mathscr{E}_\alpha=\bigcup_{\beta<\alpha}\mathscr{E}_\beta$。于是有：

% SOURCE p=55; printed=40; object=ch01-transfinite-sigma-generation; number=1.23; region=bottom
\begin{proposition}\label{proposition:ch01-transfinite-sigma-generation}
$\mathscr{M}(\mathscr{E})=\bigcup_{\alpha\in\Omega}\mathscr{E}_\alpha$，其中 $\Omega$ 为可数序数集。
\end{proposition}

% SOURCE p=55; printed=40; object=ch01-proof-transfinite-sigma-generation; region=bottom
\begin{proof}
超限归纳法表明，对所有 $\alpha\in\Omega$，有 $\mathscr{E}_\alpha\subseteq\mathscr{M}(\mathscr{E})$，因而 $\bigcup_{\alpha\in\Omega}\mathscr{E}_\alpha\subseteq\mathscr{M}(\mathscr{E})$。反向包含关系来自如下事实：$\Omega$ 中任意序列在 $\Omega$ 中都有上确界（\cref{proposition:ch00-countable-ordinals-bounded}）。若对 $j\in\mathbb{N}$ 有 $E_j\in\mathscr{E}_{\alpha_j}$，且 $\alpha=\sup\{\alpha_j\}$，则对所有 $j$ 都有 $E_j\in\mathscr{E}_\alpha$，因此 $\bigcup_1^\infty E_j\in\mathscr{E}_\beta$，其中 $\beta$ 为 $\alpha$ 的后继。
\end{proof}

% SOURCE p=56; printed=41; object=ch01-notes-sigma-cardinality; region=top
结合这一结果与\cref{proposition:ch00-continuum-operations}，可知若 $\card(\mathbb{N})\leqslant\card(\mathscr{E})\leqslant\mathfrak{c}$，则 $\card(\mathscr{M}(\mathscr{E}))=\mathfrak{c}$。（比较\cref{exercise:ch01-infinite-sigma-algebra}。）

% SOURCE p=56; printed=41; object=ch01-notes-sigma-rings; region=upper
\cref{section:ch01-measures}：有些作者倾向于取 $\sigma$-环作为测度的定义域，而不是 $\sigma$-代数（参见\cref{exercise:ch01-rings}）。原因是，在处理“非常大”的空间时，不试图测量“非常大”的集合，可以避免某些病态情形。然而，这种观点也有技术上的缺点，现在已不太流行。

% SOURCE p=56; printed=41; object=ch01-notes-extension-history; region=upper-middle
\cref{section:ch01-outer-measures}：\hyperref[theorem:ch01-caratheodory]{Carathéodory 定理}出现于他的专著\cite{folland-bib-022}。文献中曾将\cref{theorem:ch01-measure-extension}归功于 Hahn、Carathéodory 和 E. Hopf，但它最初是由 Fréchet\cite{folland-bib-054}得到的。借助\hyperref[theorem:ch01-caratheodory]{Carathéodory 定理}的证明由 Hahn\cite{folland-bib-060}和 Kolmogorov\cite{folland-bib-085}各自独立发现。

% SOURCE p=56; printed=41; object=ch01-notes-measure-construction-reference; region=middle
关于从更原始的数据构造测度这一问题的深入研究，参见 König\cite{folland-bib-086}。

% SOURCE p=56; printed=41; object=ch01-notes-lebesgue-definition; region=middle
\cref{section:ch01-borel-measures-real-line}：Lebesgue 最初和我们一样，用区间的可数覆盖来定义集合 $E\subseteq\mathbb{R}$ 的外测度 $m^*(E)$。
\begin{definition*}[Lebesgue 最初的可测性定义]
然后，他规定：有界集合 $E$ 可测，是指 $m^*(E)+m^*((a,b)\setminus E)=b-a$，其中 $(a,b)$ 是包含 $E$ 的区间；无界集合可测，是指它与任意有界区间的交可测。
\end{definition*}
Carathéodory 对可测性的刻画出现得较晚，在技术上也更便于使用。两种定义的等价性参见\cref{exercise:ch01-inner-measure}。

% SOURCE p=56; printed=41; object=ch01-notes-h-interval-complexity; region=middle-lower
思考一个 h-区间可以如何复杂地分解成 h-子区间的不交并，就能确信\cref{proposition:ch01-stieltjes-premeasure}中格外繁琐的证明是必要的。在任何这样的分解中，子区间的右端点从右向左排序后构成一个良序集，但它可以与可数序数集的任意初始段序同构。

% SOURCE p=56; printed=41; object=ch01-notes-lebesgue-extension; region=lower
Lebesgue 测度可以延拓为严格包含 $\mathscr{L}$ 的 $\sigma$-代数上的平移不变测度；参见 Kakutani 和 Oxtoby\cite{folland-bib-081}。当然，这样的 $\sigma$-代数绝不可能包含\cref{section:ch01-introduction}中讨论的不可测集。不过，Lebesgue 测度可以延拓为 $\mathscr{P}(\mathbb{R})$ 上的平移不变有限可加测度，其二维对应物（参见\cref{section:ch02-n-dimensional-lebesgue-integral}）也可以延拓为 $\mathscr{P}(\mathbb{R}^2)$ 上在平移和旋转下不变的有限可加测度；参见 Banach\cite{folland-bib-008}。Banach–Tarski 悖论使这一结果无法推广到更高维。

% SOURCE p=56; printed=41; object=ch01-notes-solovay; region=lower
关于不可测集的存在性，Solovay\cite{folland-bib-138}证明了一个引人注目的定理，其实质是：不使用选择公理，就不可能证明 Lebesgue 不可测集的存在。（该定理的精确陈述涉及公理集合论中的一些技术问题，这里不作讨论。）从从事分析研究的数学家的观点来看，Solovay 定理的意义在于重新确认 Lebesgue 理论足以满足所有实际需要。

% SOURCE p=56; printed=41; object=ch01-notes-rudin-exercise; region=bottom
\cref{exercise:ch01-borel-intermediate-density}的一个简洁解答见 Rudin\cite{folland-bib-124}。
\end{document}
